% ============================================================
% Severity Annotation Guidelines — Finance Domain (FinanceBench)
% For inclusion in EMNLP 2026 appendix or supplementary material
% ============================================================

\section{Severity Annotation Guidelines}
\label{sec:annotation-guidelines}

\subsection{Motivation}

Standard evaluation of LLMs on financial QA benchmarks treats all errors
as equally costly (binary accuracy). In practice, the cost of an error
varies by orders of magnitude: misreporting revenue by \$2B triggers
SEC enforcement and class-action litigation, while rounding a ratio to
one fewer decimal has no material impact.

We define a four-level severity taxonomy grounded in regulatory
frameworks (SEC SAB~99, SOX), enforcement data, and empirical market
evidence. Each level is associated with a representative cost $c_k$
that reflects the order-of-magnitude financial exposure of an error
at that severity.

% ============================================================
\subsection{Severity Scale}
\label{sec:severity-scale}

\begin{table}[h]
\centering
\small
\begin{tabular}{llrl}
\toprule
Level & Label & $c_k$ & Regulatory calibration \\
\midrule
1 & Negligible & \$100    & Immaterial per SAB~99; no decision impact \\
2 & Minor      & \$1\,000  & ``Little~r'' revision; isolated correction \\
3 & Major      & \$10\,000 & SOX individual penalty; operational impact \\
4 & Critical   & \$100\,000 & ``Big~R'' restatement; SEC enforcement threshold \\
\bottomrule
\end{tabular}
\caption{Four-level severity scale for the finance domain.}
\label{tab:severity-scale}
\end{table}

% ============================================================
\subsection{Empirical Calibration}
\label{sec:empirical-calibration}

The cost levels are calibrated against the following empirical evidence:

\paragraph{Negligible (\$100).}
Corresponds to the per-record cost of a data quality error in financial
systems. The ``1-10-100'' rule \citep{labovitz1992} estimates \$1 to
prevent, \$10 to correct, and \$100 per record per year for uncorrected
errors. IBM/Ponemon's 2024 Cost of a Data Breach report measures the
average cost per compromised record at \$173 \citep{ponemon2024},
consistent with this level.

\paragraph{Minor (\$1\,000).}
Corresponds to a ``little~r'' restatement---an immaterial correction
disclosed in subsequent filings without a Form~8-K Item~4.02 filing.
``Little~r'' restatements account for 76\% of all restatements as of
2020 \citep{wilmerhale2022}. The IRS accuracy-related penalty of 20\%
on a \$5\,000 understatement (\$1\,000) falls in this range. Gartner
reports that one-third of accountants make several financial errors per
week \citep{gartner2024accounting}, with individual error remediation
costs in the hundreds-to-low-thousands range.

\paragraph{Major (\$10\,000).}
The SOX Section~906 individual penalty floor is \$10\,000 for negligent
certification of inaccurate financial statements. SEC late filing
penalties range from \$10\,000 to \$97\,523 per firm
\citep{sec2025latefiling}. The ORX consortium reports a mean operational
risk loss event in banking of approximately \$250\,000 (EUR~231\,651)
\citep{orx2023}, with the distribution's lower quartile near \$10\,000.
TransAlta's 2003 spreadsheet copy-paste error cost \$24M
\citep{prosperspark2023}.

\paragraph{Critical (\$100\,000).}
The median FINRA disciplinary fine in 2024 was \$125\,000
\citep{csglaw2024}. The SEC maximum penalty per single violation is
\$775\,000 \citep{harvard2016sec}. ``Big~R'' restatements carry an
average negative impact on net income of \$16M \citep{accountingtoday2024},
with associated stock price declines of 2--10\% \citep{gao2006}.
The GAO documented over \$100B in aggregate market capitalization loss
from restatements between 1997 and 2002 \citep{gao2003}. Securities
class action median settlements are \$14M
\citep{cornerstone2024}. JPMorgan's 2012 ``London Whale'' spreadsheet
error resulted in \$6.2B in losses \citep{prosperspark2023}.

\paragraph{Beyond the scale.}
Our four-level scale deliberately caps at \$100\,000 per error instance
to model the \emph{per-query} expected cost in an AI evaluation context.
Systemic consequences (restatements, litigation, regulatory sanctions)
arise from \emph{aggregation} of errors and are captured by the compound
loss model $S = \sum_{i=1}^{N} X_i$, where even individually minor
errors can compound into critical systemic exposure.

% ============================================================
\subsection{Annotation Rules for FinanceBench}
\label{sec:annotation-rules}

Each instance $(q_i, a_i)$ in FinanceBench is annotated along two
orthogonal axes, then mapped to a severity level via a cross-tabulation
matrix.

% ------------------------------------------------------------
\subsubsection{Axis 1: Answer Type}
\label{sec:answer-type}

The answer format is a strong signal of error magnitude.

\begin{table}[h]
\centering
\small
\begin{tabular}{llp{5.5cm}}
\toprule
Code & Pattern & Detection rule \\
\midrule
\texttt{dollar\_large}
  & \$X in M/B
  & Answer contains \$-amount AND (answer mentions ``million''/``billion''
    OR value $\geq 50$ AND question mentions ``million''/``billion''
    OR value $> 500$ without scale) \\
\texttt{dollar\_small}
  & \$X per-share
  & Answer contains \$-amount with value $< 50$
    (per-share, dividend, small absolute) \\
\texttt{percentage}
  & X\%
  & Answer contains ``\%'' or ``percent'' \\
\texttt{ratio}
  & Plain number
  & Answer matches \texttt{-?[\textbackslash d,.]+x?} \\
\texttt{yes\_no}
  & Yes/No
  & Answer starts with ``yes'' or ``no'' \\
\texttt{text}
  & Free-form
  & All other answers \\
\bottomrule
\end{tabular}
\caption{Answer type classification rules (Axis~1).}
\label{tab:answer-types}
\end{table}

\paragraph{Rationale.}
A wrong absolute dollar amount in the millions (e.g., reporting revenue
as \$11.6B instead of \$14.2B) has direct financial impact: it can
trigger incorrect trading decisions, violate loan covenants, or
misrepresent materiality thresholds. A wrong percentage or ratio
(e.g., gross margin of 32\% vs.\ 34\%) is less likely to alter a
binary decision (buy/sell/hold) but still affects analytical precision.
A wrong yes/no qualitative answer about business characteristics
(``Is the company capital-intensive?'') has the least direct
dollar-denominated impact unless it concerns a regulatory or
compliance question.

% ------------------------------------------------------------
\subsubsection{Axis 2: Metric Type}
\label{sec:metric-type}

The financial metric being queried determines the \emph{domain of
impact} if the answer is wrong.

\begin{table}[h]
\centering
\small
\begin{tabular}{llp{4.5cm}}
\toprule
Code & SEC/GAAP category & Keywords (in question) \\
\midrule
\texttt{core}
  & Income statement, balance sheet line items
  & revenue, net income, net earnings, net loss,
    total assets, total liabilities, total equity,
    gross profit, COGS, cost of sales \\
\texttt{operational}
  & Cash flow, working capital, capital structure
  & EBITDA, operating income, cash flow,
    capex, capital expenditure, working capital,
    debt, inventory, accounts receivable/payable,
    pension, restructuring, derivatives \\
\texttt{ratio}
  & Derived metrics (ratios, margins, per-share)
  & ratio, margin, growth rate, return on,
    ROE, ROA, EPS, per share, yield, turnover,
    leverage, liquidity, P/E \\
\texttt{descriptive}
  & Qualitative, non-numerical
  & products and services, key agenda,
    describe, list, which securities \\
\bottomrule
\end{tabular}
\caption{Metric type classification (Axis~2). First keyword match wins.
Ordered: \texttt{descriptive} $\to$ \texttt{core} $\to$
\texttt{operational} $\to$ \texttt{ratio} $\to$ default
\texttt{descriptive}.}
\label{tab:metric-types}
\end{table}

\paragraph{Rationale.}
Core financial statement items (revenue, net income, total assets) are
the primary inputs to valuation models (DCF, comparable analysis) and
are subject to the strictest SEC materiality standards (SAB~99). An
error on these items directly affects the ``top line'' or ``bottom line''
that investors use for decision-making. Operational metrics (EBITDA,
cash flow, debt) are secondary inputs that inform creditworthiness and
operational efficiency assessments. Ratios and derived metrics are
third-order: they are computed from core/operational items and their
errors are bounded by the underlying data. Descriptive answers carry
minimal direct financial exposure.

% ------------------------------------------------------------
\subsubsection{Cross-Tabulation Matrix}
\label{sec:severity-matrix}

The severity level is determined by the intersection of the two axes:

\begin{table}[h]
\centering
\small
\begin{tabular}{l|cccc}
\toprule
& \texttt{core} & \texttt{operational} & \texttt{ratio} & \texttt{descriptive} \\
\midrule
\texttt{dollar\_large}  & \cellcolor{red!25}critical & \cellcolor{orange!25}major   & \cellcolor{orange!25}major      & \cellcolor{orange!25}major \\
\texttt{dollar\_small}   & \cellcolor{yellow!25}minor     & \cellcolor{yellow!25}minor    & \cellcolor{yellow!25}minor      & \cellcolor{green!15}negligible \\
\texttt{percentage}      & \cellcolor{yellow!25}minor     & \cellcolor{yellow!25}minor    & \cellcolor{yellow!25}minor      & \cellcolor{yellow!25}minor \\
\texttt{ratio}           & \cellcolor{yellow!25}minor     & \cellcolor{yellow!25}minor    & \cellcolor{yellow!25}minor      & \cellcolor{yellow!25}minor \\
\texttt{yes\_no}         & \cellcolor{orange!25}major     & \cellcolor{orange!25}major    & \cellcolor{yellow!25}minor      & \cellcolor{yellow!25}minor \\
\texttt{text}            & \cellcolor{orange!25}major     & \cellcolor{yellow!25}minor    & \cellcolor{yellow!25}minor      & \cellcolor{green!15}negligible \\
\bottomrule
\end{tabular}
\caption{Severity matrix: answer type (rows) $\times$ metric type
(columns). Each cell maps to one of the four severity levels.}
\label{tab:severity-matrix}
\end{table}

\paragraph{Design principles.}
\begin{enumerate}
\item \textbf{Dollar amounts in M/B always $\geq$ major.}
    A wrong figure expressed in millions or billions has irreducible
    high impact regardless of the metric category.
\item \textbf{Core + dollar\_large = critical.}
    Revenue, net income, total assets expressed as absolute dollar amounts
    in millions/billions are subject to the most stringent regulatory
    scrutiny (SAB~99 ``Big~R'' threshold).
\item \textbf{Ratios and percentages are always minor.}
    Derived metrics---even when computed from core financial items---are
    bounded in magnitude and more easily detected by cross-checking with
    source data. This reflects the LLM validation finding that the
    distinction between ``core revenue'' and ``revenue growth rate \%''
    matters more than the underlying metric category.
\item \textbf{Yes/No on core or operational topics = major.}
    Binary assessments about financial health (e.g., ``Does the company
    have adequate liquidity?'') can directly influence investment
    decisions.
\item \textbf{Text + descriptive = negligible.}
    Qualitative answers about company descriptions, filing metadata, or
    product listings have minimal direct financial exposure.
\end{enumerate}

% ============================================================
\subsection{Validation Protocol}
\label{sec:validation}

\paragraph{Step 1: Rule-based annotation.}
All 150 FinanceBench instances are annotated using the deterministic
rules above. Each annotation records the answer type, metric type, and
resulting severity, ensuring full traceability.

\paragraph{Step 2: LLM validation.}
A strong LLM (Claude Sonnet 4) independently annotates each instance
using the rubric in Table~\ref{tab:severity-scale} and a one-shot
prompt. The LLM provides a severity label and a one-sentence
justification.

\paragraph{Step 3: Disagreement resolution.}
Instances where the rule-based and LLM annotations disagree are
reviewed manually. We report the inter-annotator agreement (Cohen's
$\kappa$) between the rule-based system and the LLM.

\paragraph{Validation results.}
Table~\ref{tab:llm-agreement} reports the inter-annotator agreement
between the rule-based system and Claude Sonnet~4 on the 150
FinanceBench instances. After one round of rule refinement informed by
disagreement analysis, the system achieves moderate agreement
($\kappa = 0.46$, exact agreement 65.3\%, adjacent agreement 94.0\%).
The 52 remaining disagreements are balanced (25 LLM more severe, 27
LLM less severe) with no dominant axis, indicating no systematic bias.

\begin{table}[h]
\centering
\small
\begin{tabular}{lr}
\toprule
Metric & Value \\
\midrule
Cohen's $\kappa$ & 0.460 \\
Exact agreement & 65.3\% \\
Adjacent agreement ($\pm 1$ level) & 94.0\% \\
Disagreements & 52 / 150 \\
\quad LLM more severe & 25 \\
\quad LLM less severe & 27 \\
\bottomrule
\end{tabular}
\caption{Inter-annotator agreement: rule-based system vs.\ Claude
Sonnet~4 on FinanceBench (150 instances).}
\label{tab:llm-agreement}
\end{table}

\paragraph{Step 4: Sensitivity analysis.}
We verify that the model rankings by $\mathbb{E}[S]$ are stable under
reassignment of borderline instances ($\pm 1$ severity level) using the
sensitivity analysis module of \texttt{severity-eval}.

% ============================================================
\subsection{Empirical Evidence Summary}
\label{sec:evidence-summary}

Table~\ref{tab:evidence-summary} compiles the regulatory and market
evidence supporting the cost calibration.

\begin{table*}[t]
\centering
\small
\begin{tabular}{p{2.8cm}p{4.5cm}rp{3.5cm}}
\toprule
Source & Metric & Amount & Calibration \\
\midrule
\multicolumn{4}{l}{\textit{Per-record / per-error costs}} \\
\quad 1-10-100 Rule & Cost of uncorrected data error & \$100/record/yr & Negligible \\
\quad IBM/Ponemon 2024 & Cost per breached record & \$173 & Negligible \\
\quad IRS accuracy penalty & 20\% of understatement & \$1\,000 (on \$5k) & Minor \\
\midrule
\multicolumn{4}{l}{\textit{Regulatory penalties}} \\
\quad SOX \S906 (individual) & Negligent certification & \$10\,000 & Major \\
\quad SEC late filing & Per-firm penalty & \$10k--97k & Major \\
\quad FINRA median fine & Disciplinary action (2024) & \$125\,000 & Critical \\
\quad SEC max per violation & Single act/omission & \$775\,000 & $>$ Critical \\
\quad SOX \S906 (willful) & Willful certification & \$5\,000\,000 & $\gg$ Critical \\
\midrule
\multicolumn{4}{l}{\textit{Market impact}} \\
\quad ``Little r'' restatement & Immaterial correction & Limited & Minor--Major \\
\quad ``Big R'' restatement & Avg.\ net income impact & --\$16M & $\gg$ Critical \\
\quad Stock price (restatement) & Market-adjusted decline & --2\% to --10\% & $\gg$ Critical \\
\quad GAO 1997--2002 & Aggregate cap.\ loss & \$100B+ & Systemic \\
\midrule
\multicolumn{4}{l}{\textit{Litigation}} \\
\quad Class action (median) & Securities settlement 2024 & \$14M & $\gg$ Critical \\
\quad Class action (mean) & Securities settlement 2024 & \$42.4M & $\gg$ Critical \\
\midrule
\multicolumn{4}{l}{\textit{Operational risk}} \\
\quad ORX (banking) & Mean loss event 2023 & EUR~232k & Critical \\
\quad Spreadsheet errors & JPMorgan ``London Whale'' & \$6.2B & Catastrophic \\
\quad Spreadsheet errors & TransAlta copy-paste & \$24M & Catastrophic \\
\midrule
\multicolumn{4}{l}{\textit{Data quality (macro)}} \\
\quad Gartner & Avg.\ annual cost, poor data & \$12.9--15M/firm & Systemic \\
\quad IBM & U.S.\ annual cost, poor data & \$3.1T & Systemic \\
\quad PCAOB & Audit deficiency rate (2022) & 40\% & Prevalence \\
\quad BlackLine & CFOs not trusting data & 40\% & Prevalence \\
\bottomrule
\end{tabular}
\caption{Empirical evidence calibrating the four-level severity scale.
Amounts exceeding the \$100k ``critical'' threshold illustrate that the
per-query scale captures the \emph{unit cost} of an error; systemic
exposure is modeled through the compound loss $S = \sum X_i$.}
\label{tab:evidence-summary}
\end{table*}

% ============================================================
\subsection{Distribution on FinanceBench}
\label{sec:distribution}

Applying the annotation rules to all 150 FinanceBench instances yields:

\begin{table}[h]
\centering
\small
\begin{tabular}{lrrl}
\toprule
Level & $n$ & \% & Typical questions \\
\midrule
Negligible & 15 & 10.0\% & Company descriptions, filing metadata \\
Minor      & 86 & 57.3\% & Ratios, margins, yes/no, qualitative \\
Major      & 44 & 29.3\% & Operational \$M/B, financial health (y/n) \\
Critical   &  5 &  3.3\% & Core financials in \$M/B (revenue, net income) \\
\bottomrule
\end{tabular}
\caption{Severity distribution on FinanceBench (150 instances).}
\label{tab:financebench-distribution}
\end{table}

This distribution is consistent with the structure of the dataset:
FinanceBench contains a balanced mix of question types
(\texttt{metrics-generated}, \texttt{domain-relevant},
\texttt{novel-generated}; 50 each), and the severity distribution
reflects the natural skewness of financial risk---few questions concern
core financial statement items with catastrophic error potential, while
the majority involve derived metrics or qualitative assessments.
